\documentclass[11pt,a4paper]{article}
\usepackage[T1]{fontenc}
\usepackage[utf8]{inputenc}
\usepackage{amsmath,amssymb,amsthm}
\usepackage[margin=1in]{geometry}
\usepackage[colorlinks=true,linkcolor=blue,urlcolor=blue]{hyperref}
\theoremstyle{plain}
\newtheorem{theorem}{Theorem}
\newtheorem{lemma}[theorem]{Lemma}
\newtheorem{proposition}[theorem]{Proposition}
\newtheorem{corollary}[theorem]{Corollary}
\theoremstyle{definition}
\newtheorem{remark}[theorem]{Remark}

\title{The empirical recurrence for OEIS A229388, proved by transfer matrix}
\author{Adrian Perez Fontelles\\ \small Independent researcher}
\date{2 September 2026}
\begin{document}
\maketitle

\begin{abstract}
OEIS A229388 counts the distinct arrays obtained by applying a fixed local operation to every cell
of an array over a bounded alphabet, and carries an empirical recurrence of order $6$
contributed by R.~H.~Hardin. It is true, and it is decidable rather than empirical. What is
counted is the IMAGE of a map, not its domain, so it is not a walk count in the obvious graph;
determinising the machine that emits the derived array row by row turns it into one, on
$S=221$ states, after which the conjectured recurrence is settled by a finite exact
computation.
\end{abstract}

\noindent\small 2020 Mathematics Subject Classification. 05A15, 68Q45, 15B36.\normalsize

\section{The conjecture}

OEIS A229388 is ``Number of nX4 0..3 arrays of the sum of the corresponding element, the element to the east and the element to the south in a larger (n+1)X5 0..1 array.''. It has offset $1$ and begins
\[
164, 8592, 319340, 10624396, 345788172, 11156280332, 358453456908, 11493734735884,\ \dots
\]
The entry states, as a conjecture contributed by its author and never marked settled:
\begin{quote}
Empirical: a(n) = 49*a(n-1) -560*a(n-2) +544*a(n-3) -1568*a(n-4) +17920*a(n-5) -16384*a(n-6) for n>9
\end{quote}
As of the ``Last modified'' line on the live entry (Oct 04 2025 00:11:00, revision 8) this is still
recorded as empirical, and nothing on the entry records it as proved.

\section{What is being counted}

Let $R$ be the set of rows of length $5$ over $\{0,\dots,1\}$.
For a larger array $A=(r_1,\dots,r_{n+1})$ of $n+1$ such rows let $\Phi(A)$ be the
$n\times4$ array with
\[
\Phi(A)(i,j)\;=\;\operatorname{sum}\bigl(A(i,j),\,A(i,j+1),\,A(i+1,j)\bigr),
\]
the element, the one to its east and the one to its south. The entry counts
\[
a(n)\;=\;\bigl|\{\Phi(A)\}\bigr| ,
\]
the size of an IMAGE. That is the whole difficulty: two different larger arrays with the same
derived array must be counted once, and a walk count on the larger arrays would count them
twice.

Row $i$ of $\Phi(A)$ reads rows $i$ and $i+1$ of $A$ and nothing else, so
writing $\beta(r,x)$ for the derived row produced by two consecutive rows,
\[
\Phi(r_1,\dots,r_{n+1})=\bigl(\beta(r_1,r_2),\ \dots,\ \beta(r_n,r_{n+1})\bigr),
\]
and there is no row emitted with nothing below it.

\section{The image is a walk count after determinisation}

For a set $D$ of rows and a derived row $b$ put
\[
\delta(D,b)=\{x\in R:\ r\in D,\ \beta(r,x)=b\},
\]
and let $D_0=R$.

\begin{lemma}\label{lem:det}
For every $t\ge0$, $D_t=\delta(\dots\delta(D_0,b_1)\dots,b_t)$ is exactly the set of rows
$r_{t+1}$ completing a prefix $r_1,\dots,r_{t+1}$ whose derived rows are $b_1,\dots,b_t$.
Consequently $a(n)$ is the number of sequences $(b_1,\dots,b_n)$ with $D_n\neq\emptyset$.
\end{lemma}

\begin{proof}
Induction on $t$: $D_0=R$ is the set of possible first rows, and $\delta(D_t,b)$ collects the
rows $r_{t+2}$ for which some consistent prefix ends in a row $r_{t+1}$ with
$\beta(r_{t+1},r_{t+2})=b$. A derived array of $n$ rows is realisable exactly when the
corresponding $D_n$ is nonempty, and distinct sequences are distinct derived arrays.
\end{proof}

So with $M$ the adjacency matrix of the digraph whose vertices are the reachable nonempty sets
and which carries one edge $D\to\delta(D,b)$ for each $b$ with $\delta(D,b)\neq\emptyset$ ---
several $b$ may lead to the same set, and then the edge is repeated, which is exactly right,
since the derived arrays differ ---
\[
a(n)\;=\;\iota^{\!\top}M^{\,n}\tau,\qquad \iota=e_{D_0},\quad \tau=(1,\dots,1)^{\!\top} .
\]
There are $S=221$ reachable sets, generated from $D_0$ outward, and $a$ is $C$-finite.


\section{The proof}

\begin{lemma}\label{lem:crit}
Let $q(t)=t^{6}-\sum_i c_it^{6-i}$ be the characteristic polynomial of the
conjectured recurrence. The residual $a(n)-\sum_ic_ia(n-i)$ equals
$\iota^{\!\top}M^{\,j}q(M)\tau$ for the corresponding walk index $j$, so the recurrence
holds from some point on if and only if $\iota^{\!\top}M^{j}q(M)\tau=0$ for all large $j$,
and it is enough to examine $j<S$.
\end{lemma}

\begin{proof}
Factor $M^{j}$ out of $M^{j+6}-\sum_ic_iM^{j+6-i}$. For the bound, $M^{S}$
is an integer combination of $I,M,\dots,M^{S-1}$ by Cayley--Hamilton, so the numbers
$u_j=\iota^{\!\top}M^{j}q(M)\tau$ obey the monic recurrence given by the characteristic
polynomial of $M$; $S$ consecutive zeros force every later one, and the last nonzero $u_j$
pins the threshold exactly.
\end{proof}

\begin{theorem}
\[
a(n)\;=\;49\,a(n-1) - 560\,a(n-2) + 544\,a(n-3) - 1568\,a(n-4) + 17920\,a(n-5) - 16384\,a(n-6)
\]
for every $n>9$.
\end{theorem}

\begin{proof}
The vector $w=q(M)\tau$ was formed by $6$ matrix--vector products in exact integer
arithmetic and $\iota^{\!\top}M^{j}w$ evaluated for $j=0,1,\dots$, again exactly; those
integers vanish from the index corresponding to $n=9+1$ onwards and the run continued
until the working vector was identically zero, which settles every larger $j$ at once. The bound is exact: at $n=9$ the recurrence fails on the entry's own published terms, so no larger range holds.
\end{proof}

\section{Verification}

Three checks, each independent of the algebra above.

First, the model reproduces all $14$ terms the entry publishes, exactly, in integer
arithmetic. This is what ties the model to the entry: the machine is built by reading the
entry's English, and a misreading gives different counts.

Second, the conjectured recurrence was evaluated directly on those published terms, with no
matrices involved, and holds wherever the proved range and the data overlap.

Third, the annihilation test of Lemma~\ref{lem:crit} was carried out in exact integer
arithmetic on the whole vector rather than on a sampled prefix.

\begin{thebibliography}{9}
\bibitem{oeis} The OEIS Foundation, \emph{The On-Line Encyclopedia of Integer
Sequences}, \url{https://oeis.org}, 2026. Sequence A229388.
\bibitem{hu} J.~E.~Hopcroft, R.~Motwani and J.~D.~Ullman, \emph{Introduction to Automata
Theory, Languages, and Computation}, 3rd ed., Addison--Wesley, 2006. (The subset
construction, Section 2.3.)
\bibitem{stanley} R.~P.~Stanley, \emph{Enumerative Combinatorics}, Volume 1, 2nd ed.,
Cambridge University Press, 2011. (Transfer-matrix method, Section 4.7.)
\bibitem{horn} R.~A.~Horn and C.~R.~Johnson, \emph{Matrix Analysis}, 2nd ed., Cambridge
University Press, 2013.
\end{thebibliography}

\end{document}
